% eito para a contestação
% Cuusto de disponibilidade
% SANTO ANTÔNIO DE JESUS

\newcommand{\tipoRequerimento}{Contestação do marco temporal de devolução e revisão do período faturado}
\newcommand{\tese}{CONTESTAÇÃO}

\section{DOS FATOS E DOS DIREITOS}

\begin{enumerate}[resume]
    \item A \nomeConcessionaria{} julgou procedente a contestação apresentada quanto à Unidade Consumidora \unidadeConsumidora{}, reconhecendo o descabimento da cobrança do custo de disponibilidade para a classe de Iluminação Pública. Contudo, fundamentou indevidamente a limitação temporal do período de devolução ao marco de 31/03/2022, sob o argumento de que tal regra teria passado a vigorar apenas com a entrada em eficácia do art. 668 da REN nº 1.000/2021.
    
    \Citacao{Com a entrada em vigor da Resolução Normativa nº 1.000/2021, houve a ampliação das regras aplicáveis a toda classe de iluminação pública. Conforme estabelecido no artigo 668 da referida resolução, essas disposições passaram a ter aplicação obrigatória a partir de 31/03/2022. Dessa forma, a devolução de valores será realizada considerando como marco inicial as cobranças efetuadas a partir da ampliação da regra, em conformidade com a regulamentação vigente.}

    \item Ocorre que a vedação à cobrança de custo de disponibilidade para a classe de iluminação pública não é inovação da REN nº 1.000/2021, já sendo expressamente proibida pela regulação anterior, nos termos do art. 24, § 2º, da Resolução Normativa ANEEL nº 414/2010 ("Não se aplica a cobrança pelo custo de disponibilidade definida no art. 98 no faturamento individual de um ponto de iluminação pública"), dispositivo este sucedido e mantido pelo art. 290, § 1º, inciso I, da REN nº 1.000/2021.

    \item Além disso, a própria Resolução Normativa ANEEL nº 414/2010 já definia a classe iluminação pública de forma abrangente, contemplando o fornecimento destinado à iluminação de ruas, praças, avenidas, túneis, passagens subterrâneas, jardins, vias, estradas, passarelas, abrigos de usuários de transporte coletivo e logradouros de uso comum e livre acesso, inclusive monumentos, fachadas, fontes luminosas e obras de valor histórico, cultural ou ambiental localizadas em áreas públicas, desde que inexistente finalidade de propaganda, publicidade ou exploração econômica.

    \item Dessa forma, mesmo antes da vigência da REN ANEEL nº 1.000/2021, já havia previsão normativa suficiente para enquadrar como iluminação pública as unidades destinadas à iluminação de bens públicos de uso comum e livre acesso. Assim, não é correto afirmar que apenas a partir de 31/03/2022 teria surgido fundamento regulatório para afastar a cobrança do custo de disponibilidade.

    \item A Resolução Normativa ANEEL nº 1.000/2021 apenas reorganizou e atualizou a regulamentação setorial, preservando a vedação já existente quanto à cobrança do custo de disponibilidade em pontos de iluminação pública. Portanto, a distribuidora não pode utilizar a entrada em eficácia da nova resolução como marco inicial para limitar a devolução de valores cobrados indevidamente em período anterior, quando a regra proibitiva já constava da REN ANEEL nº 414/2010.

    \item Nesse sentido, a Unidade Consumidora \unidadeConsumidora{} estava classificada como iluminação pública, não sendo juridicamente adequado restringir a devolução apenas ao período posterior à REN ANEEL nº 1.000/2021.

    \item A própria procedência da contestação reconhecida pela distribuidora confirma que a cobrança do custo de disponibilidade era indevida para a unidade em questão. Desse modo, uma vez reconhecida a natureza da unidade como iluminação pública, a devolução deve alcançar todo o período em que a cobrança indevida ocorreu, observados os limites prescricionais (dez anos), e não apenas o período posterior a 31/03/2022.

    \item Conforme disposto no item 35 da Nota Técnica nº 136/2024/ANEEL, nos casos de condutas continuadas da distribuidora que resultem em faturamento incorreto e que abranjam períodos regidos tanto pela Resolução Normativa nº 414/2010 quanto pela nº 1.000/2021, a análise da irregularidade e o consequente dever de devolução devem ser efetuados de forma unificada, sem interrupção ou limitação arbitrária à data de mudança regulatória.

    \item No caso em análise, resta evidente que a cobrança indevida perdura continuadamente desde a vigência da REN nº 414/2010. A recusa da concessionária em restituir os valores cobrados em período anterior a 31/03/2022 descumpre o art. 323 da REN nº 1.000/2021, o Despacho ANEEL nº 2.006/2024 e o entendimento pacificado da Agência Reguladora, caracterizando retenção indevida de valores pagos a maior por erro exclusivo da distribuidora.
\end{enumerate}


\section{DOS PEDIDOS}

\begin{enumerate}[resume]
\item Ante o exposto, requer-se:
    \begin{enumerate}
        \item O afastamento do marco temporal de 31/03/2022 estabelecido pela distribuidora, reconhecendo-se que a vedação à cobrança do custo de disponibilidade para unidades classificadas como iluminação pública vigora desde a Resolução Normativa ANEEL nº 414/2010 (art. 24, § 2º), tendo sido apenas reproduzida pela Resolução Normativa ANEEL nº 1.000/2021 (art. 290, § 1º, inciso I), aplicando-se, ainda, o entendimento consolidado na Nota Técnica nº 136/2024/ANEEL;

        \item A revisão do memorial de cálculo para abranger todo o período de cobrança indevida, observado o prazo prescricional de dez anos, com a devolução, em dobro, de todos os valores indevidamente cobrados a título de custo de disponibilidade, acrescidos de correção monetária e juros legais, nos termos do art. 323 da Resolução Normativa ANEEL nº 1.000/2021, do Despacho ANEEL nº 2.006/2024 e da regulamentação aplicável.
    \end{enumerate}
\end{enumerate}